% Generated draft: requires mathematical and source review.
\documentclass[11pt]{article}
\usepackage[margin=25mm]{geometry}
\usepackage{fontspec}
\setmainfont{Latin Modern Roman}
\usepackage{amsmath,amssymb,mathtools,unicode-math}
\setmathfont{Latin Modern Math}
\usepackage{enumitem,needspace,xurl,hyperref}
\setlength{\parindent}{0pt}
\setlength{\parskip}{5pt}
\setlength{\emergencystretch}{3em}
\newcommand{\N}{\mathbb{N}}
\newcommand{\Z}{\mathbb{Z}}
\newcommand{\Q}{\mathbb{Q}}
\newcommand{\R}{\mathbb{R}}
\newcommand{\C}{\mathbb{C}}
\newcommand{\F}{\mathbb{F}}
\newcommand{\A}{\mathbb{A}}
\newcommand{\PP}{\mathbb{P}}
\newcommand{\E}{\mathbb{E}}
\newcommand{\eps}{\varepsilon}
\newcommand{\dd}{\,\mathrm{d}}
\begin{document}
\section*{UnsolvedMath ID: 701390}
\textbf{Source problem code:} catalogue entry
\subsection{Definitions and mathematical statement}
Let \(\F_2=\Z/2\Z\). Let \(T=\{0,1\}^{*}\) be the rooted binary tree of finite binary words, rooted at the empty word, with edges joining \(w\) to \(w0\) and \(w1\). For automorphisms \(u,v\) of \(T\), write \((u,v)\) for the automorphism satisfying
\[
(u,v)(0w)=0u(w),\qquad (u,v)(1w)=1v(w).
\]
Let \(1\) denote the identity automorphism and let \(a\) interchange the two subtrees by \(a(0w)=1w\) and \(a(1w)=0w\). Define \(b,c,d\) recursively by
\[
b=(a,c),\qquad c=(a,d),\qquad d=(1,b),
\]
with all four automorphisms fixing the empty word. The first Grigorchuk group is
\[
G=\langle a,b,c,d\rangle\leq\operatorname{Aut}(T).
\]
Its group algebra \(R=\F_2G\) consists of finite sums \(\sum_{g\in G}\lambda_g g\), where \(\lambda_g\in\F_2\), with coefficientwise addition and multiplication extending the group multiplication bilinearly.

Interpret the source's zeroth and first \(K\)-functors as algebraic \(K\)-theory of the unital ring \(R\). Thus \(K_0(R)\) is the Grothendieck group of the commutative monoid of isomorphism classes of finitely generated projective right \(R\)-modules under direct sum. Also, writing
\[
\operatorname{GL}(R)=\varinjlim_{n\geq1}\operatorname{GL}_n(R),
\qquad A\longmapsto\operatorname{diag}(A,1),
\]
set
\[
K_1(R)=\operatorname{GL}(R)/[\operatorname{GL}(R),\operatorname{GL}(R)].
\]
\textbf{Problem (R. Mikhailov, Kourovka Notebook 17.88).} Determine the abelian groups
\[
K_0(\F_2G)\qquad\text{and}\qquad K_1(\F_2G)
\]
up to isomorphism. The source does not specify a required form of the computation or explicitly define its \(K\)-functors; the algebraic conventions above make the intended standard interpretation explicit. The supplied registry labels this problem open; that label is not an independently verified mathematical or literature assessment.
\subsection{Short English statement}
For the first Grigorchuk group \(G\), compute the zeroth and first algebraic \(K\)-groups of its group algebra over the two-element field. The draft uses the standard projective-module definition of \(K_0\) and stable general-linear-group definition of \(K_1\).
\subsection{Sources}
\begin{itemize}
\item Local UnsolvedMath registry, permanent ID \(701390\), source problem code \texttt{catalogue entry}, title \emph{Kourovka Notebook Problem 17.88}. The supplied record attributes the problem to R. Mikhailov and gives the proposal year as 2010. Its supplied status is \texttt{open}.
\item \emph{The Kourovka Notebook (Unsolved Problems in Group Theory)}, problem 17.88, 17th issue (2010). The supplied background also identifies the 21st issue edition (2026), \texttt{arXiv:1401.0300v48}. The collection description names E. I. Khukhro and V. D. Mazurov as editors. These bibliographic details are supplied attribution and have not been independently checked.
\end{itemize}
\end{document}
